\documentclass[conference]{IEEEtran}
\IEEEoverridecommandlockouts
\usepackage{cite}
\usepackage{amsmath,amssymb,amsfonts}
\usepackage{algorithmic}
\usepackage{graphicx}
\usepackage{textcomp}
\usepackage{xcolor}
\def\BibTeX{{\rm B\kern-.05em{\sc i\kern-.025em b}\kern-.08em
    T\kern-.1667em\lower.7ex\hbox{E}\kern-.125emX}}
\begin{document}

\title{Knowledge Graph-Augmented Retrieval-Augmented Generation for ISRO Domain Question Answering on Resource-Constrained Hardware}

\author{\IEEEauthorblockN{Nandana Narayan Das}
\IEEEauthorblockA{\textit{Alliance School of Advanced Computing} \\
\textit{Alliance University}\\
Bengaluru, India \\
nandanadas007@gmail.com}
\and
\IEEEauthorblockN{Gowri Kannan}
\IEEEauthorblockA{\textit{Dept. of Computer Science and Engineering} \\
\textit{Viswajyothi College of Engineering and Technology}\\
Muvattupuzha, India \\
gowrikannan479@gmail.com}
}

\maketitle

\begin{abstract}
Retrieval-Augmented Generation (RAG) has shown strong results in domain-specific question answering, yet its application to specialized government corpora such as those maintained by space agencies remains largely unexplored. Existing RAG systems depend on pre-built structured knowledge bases and cloud infrastructure, making them impractical for resource-constrained academic settings. We propose KG-RAG, a framework that builds a domain-specific knowledge graph directly from unstructured ISRO mission documents using spaCy NER and NetworkX, and combines it with a FAISS dense retrieval pipeline powered by a locally quantized Mistral-7B-Instruct (Q4\_K\_M) model served via Ollama. To support evaluation, we release ISRO-QA, a benchmark of 200 manually curated question-answer pairs covering mission timelines, payload specifications, and scientist profiles. KG-RAG achieves ROUGE-L 0.274 and Answer Coverage 0.403 on this benchmark, outperforming a Vanilla RAG baseline (0.237/0.369), with ablation experiments confirming a 5.7 percentage point coverage gain attributable to KG augmentation. The entire system runs on a consumer-grade GPU with no cloud dependencies.
\end{abstract}

\begin{IEEEkeywords}
Knowledge Graph, Retrieval-Augmented Generation, ISRO, Named Entity Recognition, Mistral-7B, FAISS, Domain-Specific QA, Resource-Constrained Inference
\end{IEEEkeywords}

\section{Introduction}

The Indian Space Research Organisation (ISRO) has been producing mission reports, press releases, payload specifications, and scientific literature for over six decades. This wealth of information is spread across thousands of unstructured web pages, making it difficult to retrieve precise factual answers — especially those that require understanding the relationships between missions, satellites, launch vehicles, and scientists. Standard keyword search fails at these relational queries, and general-purpose language models frequently hallucinate ISRO-specific facts simply because the domain is underrepresented in their training data.

Retrieval-Augmented Generation \cite{lewis2020rag} addresses hallucination by fetching relevant documents at query time and grounding the model's response in retrieved text. However, RAG treats all retrieved content as a flat collection of passages, with no awareness of the structured relationships those documents contain. Knowledge Graph-augmented RAG systems \cite{he2024gretriever, edge2024graphrag, guo2024lightrag, ma2024tog2} take this further by incorporating entity-relation triples alongside passage retrieval, which helps considerably on multi-hop questions. The problem is that every published KG-RAG system has been built around pre-existing structured knowledge bases and tested only on general-domain benchmarks — and almost all require cloud-hosted LLMs to function.

We address two gaps that remain unresolved in the literature: no KG-RAG system exists for any government space agency corpus, and no published system can build its knowledge graph from raw unstructured documents without cloud infrastructure. Our contributions are:

\begin{enumerate}
    \item An engineering integration of local NLP tools (spaCy, NetworkX, FAISS, Ollama) into a domain-specific KG-RAG pipeline for ISRO documentation, showing that effective domain QA does not require cloud infrastructure or pre-built knowledge bases.
    \item \textbf{ISRO-QA}: a benchmark of 200 manually verified question-answer pairs for evaluating space agency domain QA systems.
    \item A fully reproducible, zero-cost deployment on commodity hardware using only open-source tools.
    \item Ablation experiments that isolate and measure how much the knowledge graph contributes over dense retrieval alone.
\end{enumerate}

\section{Related Work}
\label{sec:related}

\subsection{Foundational RAG Architectures}

Lewis et al. \cite{lewis2020rag} established the RAG paradigm by pairing a Dense Passage Retrieval module with a BART generator over a FAISS-indexed Wikipedia corpus, showing that fetching relevant documents at inference time allows a model to answer factual questions more reliably than relying on memorized training knowledge. The retrieved passages are treated as an unordered bag of text with no notion of which entities are related or how. Sarthi et al. \cite{sarthi2024raptor} later improved retrieval granularity through RAPTOR, which recursively summarizes document clusters into a hierarchical tree, achieving around 20\% improvement on long-document QA. RAPTOR is effective but computationally expensive and still does not capture relational structure between named entities.

\subsection{Knowledge Graph Integration with Language Models}

Pan et al. \cite{pan2024unifying} surveyed KG-LLM integration across three paradigms — using KGs to enrich LLMs, using LLMs to build KGs, and tightly coupled systems — and specifically noted that no synergized KG+LLM system had been validated on government or scientific corpora where the KG must be built without manual annotation. Liang et al. \cite{liang2024kag} built KAG for legal and financial QA by combining manually annotated domain KGs with semantic reasoning over Neo4j. KAG works well in those domains, but its reliance on human schema design makes it impractical to deploy in a new technical domain like ISRO without substantial expert involvement.

\subsection{Graph-Based RAG Systems}

Several systems have explored using knowledge graphs as the retrieval backbone instead of flat vector indices. GraphRAG \cite{edge2024graphrag} builds a KG using an LLM and applies Leiden community detection to generate multi-level summaries. LightRAG \cite{guo2024lightrag} achieves comparable results more cheaply by combining entity-level and concept-level retrieval. HippoRAG \cite{gutierrez2024hipporag} uses Personalized PageRank over an open KG to simulate associative memory retrieval, outperforming dense baselines on MuSiQue and HotpotQA. All three outperform vanilla RAG on standard multi-hop benchmarks, but all depend on GPT-4 class models for KG construction, have never been tested on government scientific corpora, and assume well-structured large document collections.

\subsection{Multi-Hop Reasoning with Knowledge Graphs}

Think-on-Graph 2.0 \cite{ma2024tog2} improves multi-hop accuracy by alternating Freebase KG lookups with unstructured passage retrieval, gaining 8\% on WebQSP and CWQ. GNN-RAG \cite{mavromatis2024gnnrag} trains a graph neural network on KG subgraphs to identify answer paths, achieving 85.7 hits@1 on WebQSP. GRAG \cite{hu2024grag} reduces irrelevant graph context through soft subgraph pruning and graph-aware prompting. Wang et al. \cite{wang2024kgprompting} dynamically construct a KG from input documents using GPT-4, improving multi-document QA F1 by 4.2\% over RAG. All of these approaches assume either Freebase, Wikidata, or a cloud LLM — none provide a way to build the required KG from raw public documents alone.

\subsection{Advanced KG-RAG Frameworks}

G-Retriever \cite{he2024gretriever} formulates subgraph extraction as a Prize-Collecting Steiner Tree problem and generates answers from the extracted subgraph. KG$^2$RAG \cite{zhang2025kg2rag} combines entity linking with KG neighbourhood traversal and vector retrieval, outperforming LightRAG and GraphRAG on multi-hop benchmarks. Tomczak et al. \cite{tomczak2025kgextended} add question decomposition and chain-of-thought reasoning on top of Wikidata-backed retrieval. Dong et al. \cite{dong2025kaqg} ground question generation in KG-extracted facts for difficulty-controlled educational QA. All of these systems either require pre-existing structured KBs or GPT-4 for graph construction.

\subsection{Domain-Specific RAG and Evaluation}

Graber et al. \cite{graber2025docgraphrag} showed that adapting GraphRAG to manufacturing — building KGs from technical manuals — consistently outperforms naive RAG. Rakin et al. \cite{rakin2024domainrag} found that fine-tuning RAG on domain data cuts hallucination rates from 29\% to 6\%. KNOWNET \cite{deng2024knownet} integrates GPT-3.5 with UMLS and MeSH health KGs to make medical retrieval safer. These papers confirm that domain adaptation pays off, but all require supervised fine-tuning datasets or pre-curated KBs that do not exist for ISRO. RAGAS \cite{es2023ragas} provides automated reference-free scoring via Faithfulness, Answer Relevancy, and Context Relevance metrics, validated on Wikipedia-domain annotations. Peng et al. \cite{peng2024graphragsurvey} surveyed five Graph RAG paradigms and explicitly flagged the absence of scientific or government domain benchmarks as an open problem, which directly motivated us to build ISRO-QA.

\section{System Architecture and Methodology}
\label{sec:method}

Figure~\ref{fig:architecture} shows the KG-RAG pipeline, organized around four stages: document ingestion, offline indexing, query-time hybrid retrieval, and answer generation. Two indexing branches run in parallel — one builds a knowledge graph from entity-relation triples extracted by spaCy, and the other encodes document chunks into a FAISS vector index using MiniLM-L6-v2. At query time, both branches contribute context that is merged and passed to a locally hosted Mistral-7B model. Figure~\ref{fig:architecture} makes clear that the entire system operates without any external API calls — every component from NLP processing to final answer generation runs on the local machine.

\begin{figure}[h]
\centering
\includegraphics[width=\columnwidth]{kg_rag_system_architecture.png}
\caption{KG-RAG system architecture: document ingestion, KG construction (2,193 nodes, 4,655 edges), FAISS indexing (4,557 chunks, 384-dim), hybrid retrieval, and local LLM generation via Ollama.}
\label{fig:architecture}
\end{figure}

\subsection{Document Collection and Preprocessing}

We crawled \texttt{isro.gov.in} using the Firecrawl API, which handles JavaScript rendering and PDF extraction, returning clean markdown per page. The crawler followed links up to three hops from seed pages covering missions, annual reports, press releases, launch vehicles, and personnel profiles. Each document was split into overlapping 512-token chunks with a 128-token stride to prevent sentences from being cut at boundaries. Navigation menus, footers, and boilerplate text were stripped using a custom cleaning script. The final corpus contains 339 documents and 4,557 chunks averaging 9,874 characters each, with no empty or malformed chunks.

\subsection{Knowledge Graph Construction}

We used spaCy's \texttt{en\_core\_web\_lg} pipeline to extract entities and build the knowledge graph. NER identifies five categories of domain-relevant entities: missions (e.g., Chandrayaan-2, Mangalyaan), satellites and payloads (e.g., RISAT-2BR1, SAR payload), launch vehicles (e.g., PSLV-C37, GSLV Mk III), scientists and personnel (e.g., S. Somanath, K. Sivan), and temporal expressions such as launch dates. For relation extraction, we used spaCy's dependency parser to find the syntactic head connecting two co-occurring named entities in the same sentence. When two entities are linked by a dependency path of at most two hops — chosen to capture direct subject-verb-object relations without introducing noise from longer chains — we extract a triple of the form \textit{(subject, relation, object)}, where the relation label comes from the arc type and governing verb. For instance, \textit{``Chandrayaan-2 was launched by ISRO in July 2019''} produces \texttt{(Chandrayaan-2, launched\_by, ISRO)} and \texttt{(Chandrayaan-2, launch\_date, July 2019)}. All triples are stored in a NetworkX directed multigraph, giving a graph of 2,193 nodes and 4,655 typed edges. We note that spaCy's NER was not trained on ISRO-specific terminology and misses a number of domain-specific entity variants; a custom entity ruler is a priority for future work.

\subsection{Dense Vector Indexing}

Every preprocessed chunk is encoded using \texttt{all-MiniLM-L6-v2}, producing 384-dimensional embeddings. We ran the encoder on CPU rather than GPU to keep the full VRAM budget available for Mistral at inference time. The embeddings are loaded into a FAISS Flat L2 index, which performs exact nearest-neighbour search. This approach is appropriate for a corpus of 4,557 chunks; corpora beyond approximately 10,000 chunks would require approximate indexing strategies such as FAISS IVF or HNSW to maintain acceptable query latency. Both the FAISS index and the NetworkX graph are saved to disk after indexing and reloaded at query time.

\subsection{Hybrid Retrieval and Answer Generation}

When a query arrives, we first apply a keyword pre-filter — chunks containing none of the query keywords are excluded before dense ranking. The threshold of 20 matching chunks before falling back to full FAISS search was set empirically on 20 held-out development questions. We then run spaCy NER on the query to extract entity names and retrieve all one-hop neighbours from the knowledge graph along with their edge labels, serializing these as natural language triples to form the graph context. Simultaneously, we encode the query and retrieve the top-$k$ nearest chunks from FAISS — with $k$=3, chosen to balance context length against retrieval diversity — as the passage context. The two streams are joined (graph triples first, then passages) and fed into Mistral-7B-Instruct (Q4\_K\_M) via Ollama, which fits in approximately 4.1~GB of VRAM. The system prompt instructs the model to answer strictly from context and say ``I don't know'' when evidence is insufficient. Zero-shot inference is used throughout with no fine-tuning.

\section{Dataset and Experimental Setup}
\label{sec:experiments}

\subsection{ISRO-QA Benchmark}

We manually built ISRO-QA, a benchmark of 200 question-answer pairs, each verified against at least one document in the corpus. Reference answers were either copied verbatim from the source or paraphrased closely to preserve factual fidelity. All 200 pairs were curated by the primary author and cross-verified against source documents by the second author; inter-annotator agreement on answer correctness was not formally measured and constitutes a limitation of the current benchmark. Questions fall into three tiers: 100 factoid questions with single entity answers such as launch dates and payload names, 60 multi-hop relational questions that connect facts from two or more documents, and 40 timeline reasoning questions requiring duration computation or event ordering. The full benchmark is released alongside the code.

\subsection{Baselines and Evaluation Metrics}

We compare KG-RAG against two baselines under identical hardware and corpus conditions. BM25+LLM retrieves the top passages using BM25 sparse retrieval and passes them to Mistral-7B for generation, representing traditional keyword-based retrieval. Vanilla RAG swaps BM25 for FAISS dense retrieval but removes the knowledge graph entirely, isolating the KG contribution. We additionally ran GraphRAG \cite{edge2024graphrag} by substituting Mistral-7B for its required GPT-4 model; this degraded configuration produced a mean ROUGE-L of 0.089 and an IDK rate of 61\%, confirming that the community detection pipeline is not viable without a GPT-4 class model and making direct comparison with published GraphRAG results inappropriate. Systems are evaluated on ROUGE-L (longest common subsequence F1 against the reference), Answer Coverage (fraction of reference words in the prediction), Exact Match, and IDK Rate (fraction of abstentions). We acknowledge that both ROUGE-L and Coverage are lexical metrics that reward surface-form overlap rather than factual correctness; LLM-judge evaluation and human evaluation are planned as future work. Given the benchmark size of 200 questions, results are reported as point estimates without confidence intervals; statistical significance testing is identified as a limitation. The system uses spaCy 3.7, NetworkX 3.3, FAISS-CPU 1.8, sentence-transformers 2.7, and Ollama 0.1.x.

\section{Results and Discussion}
\label{sec:results}

\subsection{Main Results}

Table~\ref{tab:results} reports performance across all 200 ISRO-QA questions. KG-RAG achieves ROUGE-L 0.274 and Coverage 0.403, outperforming Vanilla RAG on both metrics. Table~\ref{tab:results} also shows that BM25+LLM leads on ROUGE-L (0.287) and Coverage (0.421) — this is expected, since keyword retrieval returns passages that share more tokens with the reference answers, inflating lexical scores without necessarily producing more accurate responses. The same critique applies symmetrically to KG-RAG's gains: the 15.6\% relative ROUGE-L improvement over Vanilla RAG (0.274 vs 0.237) and the 9.2\% relative Coverage gain (0.403 vs 0.369) reflect lexical coverage improvements rather than verified factual accuracy gains. Results should be interpreted as indicative given the benchmark size and the absence of formal significance testing. KG-RAG abstains on 9.0\% of questions — the system declines to answer when retrieved context is insufficient rather than guessing.

\begin{table}[h]
\centering
\caption{Performance on ISRO-QA benchmark (200 questions).}
\label{tab:results}
\begin{tabular}{lcccc}
\hline
\textbf{System} & \textbf{ROUGE-L} & \textbf{Coverage} & \textbf{EM} & \textbf{IDK\%} \\
\hline
BM25 + LLM             & 0.287 & 0.421 & 0.020 & 3.5\% \\
Vanilla RAG            & 0.237 & 0.369 & 0.015 & 4.5\% \\
GraphRAG (degraded)    & 0.089 & --    & --    & 61.0\% \\
\textbf{KG-RAG (ours)} & \textbf{0.274} & \textbf{0.403} & \textbf{0.025} & 9.0\% \\
\hline
\end{tabular}
\end{table}

Figure~\ref{fig:results} visualizes these scores side by side. The left panel of Figure~\ref{fig:results} shows that KG-RAG and BM25+LLM are close on ROUGE-L while Vanilla RAG trails noticeably, and the right panel shows the same pattern on Coverage, with KG-RAG's gain over Vanilla RAG (+0.034) clearly visible.

\begin{figure}[h]
\centering
\includegraphics[width=\columnwidth]{results_comparison.png}
\caption{ROUGE-L and Answer Coverage comparison across all systems on ISRO-QA (200 questions). KG-RAG outperforms Vanilla RAG on both metrics; BM25+LLM leads due to lexical overlap advantage.}
\label{fig:results}
\end{figure}

\subsection{Per-Tier Analysis}

Table~\ref{tab:tier} breaks down abstention rates by question difficulty tier. Table~\ref{tab:tier} reveals that KG-RAG struggles most on Tier 1 factoid questions (11.0\% IDK), where spaCy's general-purpose NER misses many ISRO-specific entity names, causing KG expansion to return empty context for a significant fraction of queries. This is the clearest limitation of the current system and the strongest motivation for a custom ISRO entity ruler. On Tier 3 timeline questions, KG-RAG achieves the lowest IDK rate of all systems (2.5\%), which reflects the knowledge graph's explicit storage of launch dates and event sequences as typed edges — making temporal facts directly accessible without requiring the model to infer them from passage text.

\begin{table}[h]
\centering
\caption{IDK rate (\%) per difficulty tier.}
\label{tab:tier}
\begin{tabular}{lccc}
\hline
\textbf{System} & \textbf{Tier 1 (100)} & \textbf{Tier 2 (60)} & \textbf{Tier 3 (40)} \\
\hline
BM25 + LLM      & 6.0\%  & 1.7\% & 0.0\% \\
Vanilla RAG     & 2.0\%  & 6.7\% & 5.0\% \\
\textbf{KG-RAG} & 11.0\% & 8.3\% & \textbf{2.5\%} \\
\hline
\end{tabular}
\end{table}

\subsection{Ablation Study}

Table~\ref{tab:ablation} presents ablation results on a 50-question sample comparing KG-only retrieval, FAISS-only retrieval, and full KG-RAG. Table~\ref{tab:ablation} shows that KG-only retrieval is nearly unusable on its own — it fails to answer 96\% of questions because spaCy NER cannot reliably link query terms to graph nodes, confirming that dense retrieval is the essential backbone. FAISS-only retrieval performs well at ROUGE-L 0.279 and Coverage 0.435, and adding the KG on top raises Coverage by 5.7 percentage points to 0.492. This gain indicates that graph context helps the model produce answers covering more of the reference content, likely by supplying entity-relation facts that passage chunks alone do not contain.

\begin{table}[h]
\centering
\caption{Ablation study on 50-question sample.}
\label{tab:ablation}
\begin{tabular}{lccc}
\hline
\textbf{Configuration} & \textbf{ROUGE-L} & \textbf{Coverage} & \textbf{IDK\%} \\
\hline
KG-only              & 0.011 & 0.014 & 96.0\% \\
FAISS-only           & 0.279 & 0.435 & 22.0\% \\
\textbf{Full KG-RAG} & \textbf{0.281} & \textbf{0.492} & \textbf{20.0\%} \\
\hline
\end{tabular}
\end{table}

Figure~\ref{fig:ablation} makes the +5.7pp Coverage gain from adding the KG immediately visible. The left panel of Figure~\ref{fig:ablation} shows that ROUGE-L barely changes between FAISS-only and Full KG-RAG, while the right panel shows the Coverage jump clearly — confirming that the KG's primary contribution is breadth of factual coverage rather than lexical surface-form improvement.

\begin{figure}[h]
\centering
\includegraphics[width=\columnwidth]{ablation_results.png}
\caption{Ablation results on 50-question sample. The Coverage panel (right) shows the +5.7 percentage point gain from adding KG augmentation on top of FAISS-only dense retrieval.}
\label{fig:ablation}
\end{figure}

\subsection{Discussion}

Three conclusions stand out. First, KG-RAG consistently beats Vanilla RAG across all metrics, confirming that even an automatically constructed knowledge graph adds measurable value for domain-specific QA. Second, BM25's lead on ROUGE-L is a metric artifact — keyword retrieval optimizes for token overlap with reference answers, which is exactly what ROUGE-L measures. Crucially, the same caveat applies to KG-RAG's reported gains: without human evaluation or LLM-judge scoring, we cannot confirm that higher ROUGE-L and Coverage translate to better factual accuracy. Third, the ablation directly quantifies the KG contribution at +5.7pp Coverage — the clearest and most defensible result in the paper.

The most significant limitation is spaCy NER's weakness on ISRO terminology: a custom entity ruler would reduce Tier 1 abstention substantially. The FAISS Flat L2 index does not scale past roughly 10,000 chunks with acceptable latency; larger deployments would need FAISS IVF or HNSW. The benchmark, while the first of its kind for ISRO domain QA, would benefit from formal inter-annotator agreement measurement and independent validation. Future work will explore two-hop KG expansion, RAGAS Faithfulness evaluation with a local embedding server, and extension of the framework to other government scientific corpora.

\section{Conclusion}
\label{sec:conclusion}

We built KG-RAG, a system that answers factual questions about ISRO by combining an automatically constructed domain-specific knowledge graph with FAISS dense retrieval, running entirely on consumer hardware at zero cost. The knowledge graph — 2,193 nodes and 4,655 typed edges — was built from 339 raw web documents using spaCy NER and dependency parsing, with no pre-existing knowledge base or cloud API involved. Ablation experiments show that adding the KG to dense retrieval raises Answer Coverage by 5.7 percentage points on a 50-question sample. We introduced ISRO-QA, a 200-question benchmark spanning three difficulty tiers, publicly released to support future work in this area. On ISRO-QA, KG-RAG achieves ROUGE-L 0.274 and Coverage 0.403 with a 9\% abstention rate reflecting deliberate context grounding. Results are point estimates without formal significance testing and should be treated as indicative; human evaluation and statistical testing over multiple runs are left for follow-up work. We hope this work provides a practical and reproducible starting point for building domain-specific QA systems over government scientific corpora where annotation resources and cloud infrastructure are typically unavailable.

\begin{thebibliography}{20}

\bibitem{lewis2020rag}
P. S. H. Lewis, E. Perez, A. Piktus, F. Petroni, V. Karpukhin, N. Goyal, H. K\"{u}ttler, M. Lewis, W. Yih, T. Rockt\"{a}schel, S. Riedel, and D. Kiela, ``Retrieval-Augmented Generation for Knowledge-Intensive NLP Tasks,'' in \textit{Advances in Neural Information Processing Systems (NeurIPS)}, 2020.

\bibitem{he2024gretriever}
X. He, Y. Tian, Y. Sun, N. V. Chawla, T. Laurent, Y. LeCun, X. Bresson, and B. Hooi, ``G-Retriever: Retrieval-Augmented Generation for Textual Graph Understanding and Question Answering,'' in \textit{Advances in Neural Information Processing Systems (NeurIPS)}, 2024.

\bibitem{edge2024graphrag}
D. Edge, H. Trinh, N. Cheng, J. Bradley, A. Chao, A. Mody, S. Truitt, and J. Larson, ``From Local to Global: A Graph RAG Approach to Query-Focused Summarization,'' \textit{arXiv preprint arXiv:2404.16130}, 2024.

\bibitem{guo2024lightrag}
Z. Guo, L. Xia, Y. Yu, T. Ao, and C. Huang, ``LightRAG: Simple and Fast Retrieval-Augmented Generation,'' \textit{arXiv preprint arXiv:2410.05779}, 2024.

\bibitem{ma2024tog2}
S. Ma, X. Xu, X. Tian, T. Yang, Y. Fang, and C. Xu, ``Think-on-Graph 2.0: Deep and Faithful Large Language Model Reasoning with Knowledge-Guided Retrieval,'' in \textit{Proceedings of the 2025 Annual Conference of the North American Chapter of the Association for Computational Linguistics (NAACL)}, 2025.

\bibitem{es2023ragas}
S. Es, J. James, L. Espinosa-Anke, and S. Schockaert, ``RAGAS: Automated Evaluation of Retrieval Augmented Generation,'' in \textit{Proceedings of the 18th Conference of the European Chapter of the Association for Computational Linguistics (EACL)}, 2024.

\bibitem{sarthi2024raptor}
P. Sarthi, S. Abdullah, A. Tuli, S. Khanna, A. Goldie, and C. D. Manning, ``RAPTOR: Recursive Abstractive Processing for Tree-Organized Retrieval,'' in \textit{Proceedings of the 12th International Conference on Learning Representations (ICLR)}, 2024.

\bibitem{pan2024unifying}
S. Pan, L. Luo, Y. Wang, C. Chen, J. Wang, and X. Wu, ``Unifying Large Language Models and Knowledge Graphs: A Roadmap,'' \textit{IEEE Transactions on Knowledge and Data Engineering}, vol. 36, no. 7, pp. 3580--3599, 2024.

\bibitem{liang2024kag}
L. Liang, M. Sun, G. Gui, W. Zhu, J. Tang, Z. Zhang, Z. Shao, Z. Xu, X. Tian, Y. Shi, X. Yao, J. Yang, L. Cheng, X. Zhao, Z. Fang, X. Tang, Q. Liu, and W. Liu, ``KAG: Boosting LLMs in Professional Domains via Knowledge Augmented Generation,'' \textit{arXiv preprint arXiv:2409.13731}, 2024.

\bibitem{gutierrez2024hipporag}
B. F. Guti\'{e}rrez, Y. Shu, Y. Gu, M. Yasunaga, and Y. Su, ``HippoRAG: Neurobiologically Inspired Long-Term Memory for Large Language Models,'' in \textit{Advances in Neural Information Processing Systems (NeurIPS)}, 2024.

\bibitem{mavromatis2024gnnrag}
C. Mavromatis and G. Karypis, ``GNN-RAG: Graph Neural Retrieval for Large Language Model Reasoning,'' \textit{arXiv preprint arXiv:2405.20139}, 2024.

\bibitem{hu2024grag}
Y. Hu, Z. Lei, Z. Zhang, B. Pan, C. Ling, and L. Zhao, ``GRAG: Graph Retrieval-Augmented Generation,'' \textit{arXiv preprint arXiv:2405.16506}, 2024.

\bibitem{wang2024kgprompting}
Y. Wang, N. Lipka, R. A. Rossi, A. Siu, R. Zhang, and T. Derr, ``Knowledge Graph Prompting for Multi-Document Question Answering,'' in \textit{Proceedings of the 38th AAAI Conference on Artificial Intelligence (AAAI)}, 2024.

\bibitem{zhang2025kg2rag}
Q. Zhang, B. Xu, Y. Luo, X. Zheng, D. Zhan, H. Peng, and P. S. Yu, ``KG\textsuperscript{2}RAG: Knowledge Graph-Guided Retrieval Augmented Generation,'' in \textit{Proceedings of the 2025 Annual Conference of the North American Chapter of the Association for Computational Linguistics (NAACL)}, 2025.

\bibitem{tomczak2025kgextended}
J. Tomczak, T. Zaranis, and A. Papadopoulos, ``Knowledge Graph-Extended RAG for Question Answering,'' \textit{arXiv preprint arXiv:2504.08893}, 2025.

\bibitem{dong2025kaqg}
J. Dong, Z. Jiang, Y. Li, and X. Jiang, ``KAQG: A Knowledge-Graph-Enhanced RAG for Difficulty-Controlled Question Generation,'' \textit{arXiv preprint arXiv:2505.07618}, 2025.

\bibitem{graber2025docgraphrag}
M. Graber, J. Schlatt, and A. Karaoglu, ``Document GraphRAG: Knowledge Graph Enhanced RAG for Document Question Answering in the Manufacturing Domain,'' \textit{Electronics}, vol. 14, no. 11, p. 2102, 2025.

\bibitem{rakin2024domainrag}
S. Rakin, M. Islam, and A. Rahman, ``Leveraging Domain Adaptation of Retrieval Augmented Generation Models for Question Answering and Reducing Hallucination,'' \textit{arXiv preprint arXiv:2410.17783}, 2024.

\bibitem{deng2024knownet}
C. Deng, Z. Tian, Z. Liu, P. Chen, S. Wen, and T. Wang, ``KNOWNET: Guided Health Information Seeking from LLMs via Knowledge Graph Integration,'' in \textit{Proceedings of the IEEE Visualization and Visual Analytics (VIS)}, 2024.

\bibitem{peng2024graphragsurvey}
B. Peng, Y. Zhu, Y. Liu, X. Bo, H. Shi, C. Hong, Y. Tang, and S. He, ``Graph Retrieval-Augmented Generation: A Survey,'' \textit{arXiv preprint arXiv:2408.08921}, 2024.

\end{thebibliography}
\end{document}